\begin{table}[!htbp]
\centering
\caption{Within-cell associations between sex-specific employment composition and automation margins.}
\label{tab:appendix_ilostat_gender_fe_baseline}
\small
\begin{tabular}{lcccc}
\toprule
 & (1) & (2) & (3) & (4) \\
 & Occupation sub. & Occupation aug. & Industry sub. & Industry aug. \\
\midrule
Exposure margin (10 pp) & -0.338*** & 0.140 & -0.227*** & 0.015 \\
 & (0.123) & (0.429) & (0.042) & (0.076) \\
$t$-statistic & -2.74 & 0.33 & -5.46 & 0.20 \\
$p$-value & 0.007 & 0.744 & $<0.001$ & 0.842 \\
\midrule
Observations & 3,042 & 3,042 & 4,760 & 4,760 \\
Countries & 88 & 88 & 72 & 72 \\
Cells & 42 & 42 & 88 & 88 \\
Country fixed effects & Yes & Yes & Yes & Yes \\
Cell fixed effects & Yes & Yes & Yes & Yes \\
\bottomrule
\end{tabular}
\caption*{\scriptsize Notes: The outcome is the percentage-point difference between the share of female employment and the share of male employment located in a country--cell. Female and male shares are normalised over cells observed for both sexes within each country. Each column is a separate regression with country and exact two-digit occupation or industry cell fixed effects. Exposure margins are scaled so that coefficients correspond to a 10 percentage point increase. Negative coefficients indicate that a cell accounts for a relatively larger share of male employment; they do not imply that men constitute a majority of workers in that cell. The sample is restricted to countries classified in the four World Bank income groups, matching Figure~\ref{fig:gender_informality_margin_applications}. Standard errors, clustered by country, are in parentheses. Tests are two-sided. $^{*}p<0.10$, $^{**}p<0.05$, $^{***}p<0.01$.}
\end{table}
